\modulechapter{8}{Ingestion và các pattern pipeline}
\moduleinfo{Trung cấp}{20 giờ}{Module 5, 7}

Ingestion là bước đưa dữ liệu từ hệ thống nguồn vào tầng raw của lakehouse. Bài toán khó nhất không phải đọc dữ liệu một lần, mà đọc lặp lại mỗi ngày sao cho không sót, không trùng và chạy lại được khi lỗi. Module này đặt tên cho các pattern chuẩn: nạp toàn bộ và nạp tăng dần, cursor và high-water mark, idempotency, dữ liệu tới muộn, backfill, CDC và schema drift. Học xong, bạn xây được một pipeline tăng dần từ PostgreSQL vào Iceberg mà chạy lại bao nhiêu lần cũng cho cùng kết quả.

\section{Bạn sẽ học được}
\begin{itemize}
  \item Chọn giữa nạp toàn bộ và nạp tăng dần, và gọi đúng tên các sync mode thường gặp.
  \item Thiết kế cursor hoặc high-water mark, và chỉ ra những trường hợp nó bỏ sót thay đổi.
  \item Viết bước ghi idempotent bằng \code{MERGE} hoặc overwrite partition, và lưu tiến độ đúng thời điểm.
  \item Xử lý dữ liệu tới muộn bằng cửa sổ nhìn lại, và lập kế hoạch backfill.
  \item Giải thích CDC dựa trên log khác truy vấn theo \code{updated\_at} ở đâu, và đọc được một sự kiện thay đổi của Debezium.
  \item Phân loại thay đổi schema và chọn phản ứng phù hợp ở từng tầng.
\end{itemize}

\section{Nội dung chính}
\begin{enumerate}
  \item \textbf{Nguồn và tần suất.} Nguồn có thể là cơ sở dữ liệu OLTP, API, file hay luồng sự kiện; bạn kéo (pull) theo lịch hoặc nhận đẩy (push). Batch và streaming là hai đầu của một dải tần suất; batch chạy mỗi vài phút đã gần với micro-batch.
  \item \textbf{Nạp toàn bộ và nạp tăng dần.} Airbyte đặt tên năm sync mode: Incremental Append, Incremental Append + Deduped, Full Refresh Append, Full Refresh Overwrite, Full Refresh Overwrite + Deduped; các mode Deduped khử trùng theo primary key (\href{https://docs.airbyte.com/platform/using-airbyte/core-concepts/sync-modes}{Airbyte Sync Modes}). Dù dùng công cụ nào, đây là bộ từ vựng chung.
  \item \textbf{Cursor và high-water mark.} Nạp tăng dần lưu giá trị lớn nhất đã nạp của một cột (cursor) rồi lần sau chỉ đọc phần lớn hơn. dlt lưu giá trị này trong state và đưa ra dưới dạng \code{start\_value} ở lần chạy sau; biên dưới mặc định là bao gồm, nên dòng ở đúng giá trị cũ được đọc lại và dlt khử trùng chúng (\href{https://dlthub.com/docs/general-usage/incremental/cursor}{dlt: cursor-based incremental loading}). Đừng nhầm high-water mark của batch với watermark của streaming, thứ ước lượng tiến độ theo event time.
  \item \textbf{Cursor bỏ sót gì.} Cursor trên cột thời điểm tạo (ví dụ \code{order\_purchase\_timestamp}) chỉ bắt dòng mới, không bắt dòng bị sửa. Cursor trên \code{updated\_at} bắt được sửa nếu nguồn luôn cập nhật cột này, nhưng không bao giờ thấy dòng bị xoá cứng. Chốt cận trên của mỗi lần chạy (ví dụ thời điểm bắt đầu run) để hai lần chạy không chồng nhau một cách khó đoán.
  \item \textbf{Lưu tiến độ sau khi dữ liệu đã commit.} Nếu ghi cursor trước rồi mới ghi dữ liệu, một lần lỗi giữa chừng làm mất dữ liệu. Ghi dữ liệu trước, cursor sau: lỗi giữa chừng chỉ làm đọc lại một phần, và đích idempotent biến việc đọc lại thành vô hại. Đây là at-least-once cộng ghi idempotent.
  \item \textbf{Idempotency.} Task idempotent cho cùng kết quả mỗi lần chạy lại. Docs Airflow khuyên coi task như transaction, không dùng \code{INSERT} khi chạy lại mà dùng UPSERT, đọc và ghi đúng một partition xác định, và không dùng \code{datetime.now()} cho phép tính quan trọng (\href{https://airflow.apache.org/docs/apache-airflow/stable/best-practices.html}{Airflow Best Practices}). Trên Iceberg, hai cách chuẩn là \code{MERGE INTO} theo khoá và overwrite partition động (\href{https://iceberg.apache.org/docs/latest/spark-writes/}{Iceberg Spark Writes}).
  \item \textbf{Dữ liệu tới muộn.} Dòng có thể tới sau khi cursor đã vượt qua nó. Cách phổ biến là cửa sổ nhìn lại: mỗi lần đọc lại một khoảng trước cursor rồi \code{MERGE}. dlt có tuỳ chọn \code{lag} cho việc này, tính bằng giây với cursor kiểu datetime, bằng ngày với cursor kiểu date (\href{https://dlthub.com/docs/general-usage/incremental/lag}{dlt: lag}).
  \item \textbf{Backfill và xử lý lại.} Backfill là chạy lại pipeline cho một khoảng thời gian trong quá khứ, ví dụ sau khi sửa logic. Nó chỉ an toàn khi mỗi lần chạy ghi đè đúng phần của nó. Airflow 3 có backfill do scheduler quản lý (\href{https://airflow.apache.org/docs/apache-airflow/stable/core-concepts/backfill.html}{Airflow Backfill}); Module 9 sẽ dùng nó.
  \item \textbf{CDC dựa trên log.} Thay vì truy vấn bảng, CDC đọc log của cơ sở dữ liệu. Connector PostgreSQL của Debezium dùng logical decoding với plugin \code{pgoutput} (có sẵn từ PostgreSQL 10) qua một replication slot; lần đầu nó chụp snapshot, mỗi dòng snapshot là sự kiện \code{op=r}, rồi đọc tiếp từ WAL (\href{https://debezium.io/documentation/reference/stable/connectors/postgresql.html}{Debezium PostgreSQL connector}, \href{https://www.postgresql.org/docs/18/logicaldecoding.html}{PostgreSQL: Logical Decoding}).
  \item \textbf{Cấu trúc sự kiện thay đổi.} Payload có \code{before}, \code{after}, \code{source}, \code{op} và \code{ts\_ms}; \code{op} là \code{c} (tạo), \code{u} (sửa), \code{d} (xoá), \code{r} (đọc snapshot), \code{t} (truncate). Nội dung \code{before} phụ thuộc \code{REPLICA IDENTITY} của bảng: mặc định chỉ có khoá chính, \code{FULL} thì có mọi cột. Sau sự kiện xoá, connector gửi thêm một tombstone (khoá giữ nguyên, value null) để Kafka log compaction dọn được khoá đó (\href{https://debezium.io/documentation/reference/stable/connectors/postgresql.html}{Debezium PostgreSQL connector}).
  \item \textbf{Kafka Connect.} Kafka Connect là công cụ truyền dữ liệu giữa Kafka và hệ thống khác một cách có thể mở rộng và tin cậy; nó chạy ở chế độ standalone hoặc distributed và tự quản lý việc commit offset cho connector (\href{https://kafka.apache.org/44/kafka-connect/overview/}{Kafka Connect Overview}). Debezium chạy như một source connector trên Kafka Connect (\href{https://debezium.io/documentation/reference/stable/architecture.html}{Debezium Architecture}).
  \item \textbf{Schema drift.} Thêm cột là thay đổi tương thích; đổi kiểu, đổi tên hay xoá cột có thể phá vỡ tầng sau. Iceberg hỗ trợ schema evolution an toàn vì cột được theo dõi bằng ID chứ không bằng tên (\href{https://iceberg.apache.org/docs/latest/evolution/}{Iceberg Evolution}). Ở tầng raw thường cho phép evolve; ở staging và marts thì nên cố định và kiểm bằng test hoặc contract (Module 10).
  \item \textbf{ETL và ELT.} ETL biến đổi trước khi nạp vào đích; ELT nạp nguyên bản vào tầng raw rồi biến đổi bên trong lakehouse. Kiến trúc raw, staging, marts của lộ trình này là ELT.
\end{enumerate}

\section{Bắt đầu}
\begin{itemize}
  \item \href{https://docs.airbyte.com/platform/using-airbyte/core-concepts/sync-modes}{Airbyte: Sync Modes}: đọc 15 phút để nắm từ vựng.
  \item \href{https://dlthub.com/docs/general-usage/incremental/cursor}{dlt: cursor-based incremental loading}: chạy ví dụ trong trang với destination DuckDB.
  \item \href{https://debezium.io/documentation/reference/stable/tutorial.html}{Debezium Tutorial}: chạy theo tutorial chính thức để thấy sự kiện create, update, delete.
\end{itemize}

\section{Tài liệu chính thức}
\begin{xltabular}{\textwidth}{@{}P{4.6cm} L P{1.8cm}@{}}
\toprule
\thd{tài liệu} & \thd{nội dung cần đọc} & \thd{thời lượng} \\
\midrule
\endhead
\bottomrule
\endfoot
\href{https://docs.airbyte.com/platform/using-airbyte/core-concepts/sync-modes}{Airbyte: Sync Modes} & năm sync mode và cách khử trùng theo primary key & 0,5 giờ \\
\href{https://dlthub.com/docs/general-usage/incremental/cursor}{dlt: cursor-based incremental loading}, \href{https://dlthub.com/docs/general-usage/incremental/lag}{lag} & \code{dlt.sources.incremental}, state, khử trùng ở biên, \code{row\_order}, \code{lag} & 2 giờ \\
\href{https://airflow.apache.org/docs/apache-airflow/stable/best-practices.html}{Airflow: Best Practices} & mục Creating a task: task như transaction, UPSERT, partition xác định & 0,5 giờ \\
\href{https://iceberg.apache.org/docs/latest/spark-writes/}{Iceberg: Spark Writes} & MERGE INTO, INSERT OVERWRITE động & 1 giờ \\
\href{https://spark.apache.org/docs/latest/sql-data-sources-jdbc.html}{Spark: JDBC To Other Databases} & các option \code{url}, \code{query}, \code{dbtable}, driver & 0,5 giờ \\
\href{https://debezium.io/documentation/reference/stable/architecture.html}{Debezium: Architecture} & Debezium trên Kafka Connect, Debezium Server & 0,5 giờ \\
\href{https://debezium.io/documentation/reference/stable/connectors/postgresql.html}{Debezium: PostgreSQL connector} & How the connector works, snapshot, data change events & 2 giờ \\
\href{https://kafka.apache.org/44/kafka-connect/overview/}{Kafka Connect} & Overview và User Guide: standalone, distributed, offset & 1 giờ \\
\href{https://iceberg.apache.org/docs/latest/evolution/}{Iceberg: Evolution} & schema evolution, partition evolution & 0,5 giờ \\
\href{https://airflow.apache.org/docs/apache-airflow/stable/core-concepts/backfill.html}{Airflow: Backfill} & backfill do scheduler quản lý & 0,5 giờ \\
\end{xltabular}

\section{Khoá học miễn phí}
\begin{xltabular}{\textwidth}{@{}P{4.2cm} L P{2.6cm}@{}}
\toprule
\thd{khoá} & \thd{nội dung} & \thd{miễn phí} \\
\midrule
\endhead
\bottomrule
\endfoot
\href{https://github.com/DataTalksClub/data-engineering-zoomcamp}{DE Zoomcamp (DataTalksClub)} & khoá data engineering mở, có workshop nạp dữ liệu với dlt & toàn bộ tài liệu trên GitHub \\
\href{https://developer.confluent.io/courses/}{Confluent Developer: courses} & các khoá về Kafka, Kafka Connect và Schema Registry; học sau khi làm lab Debezium & khoá tự học trên trang \\
\end{xltabular}

\section{Lab thực hành}

\subsection{Lab 8.1: dựng nguồn và nạp toàn bộ}
\begin{enumerate}
  \item Nạp \code{olist\_orders\_dataset.csv} vào PostgreSQL thành bảng \code{src.orders}, thêm cột \code{updated\_at timestamptz} và đặt bằng \code{order\_purchase\_timestamp} cho mọi dòng hiện có.
  \item Để giả lập nguồn đang sống, chỉ giữ lại các đơn tới hết tháng 06/2017; giữ phần còn lại trong một bảng chờ \code{src.orders\_future}.
  \item Từ Spark, đọc \code{src.orders} qua JDBC (driver PostgreSQL phải có trên classpath của Spark server) và ghi nguyên bản vào \code{lake.raw.orders} bằng overwrite.
  \item Chạy lại bước 3 hai lần; so số dòng giữa nguồn và đích.
\end{enumerate}
\hoanthanh{số dòng của \code{lake.raw.orders} bằng \code{src.orders} sau mỗi lần chạy, và bạn giải thích được vì sao nạp toàn bộ bằng overwrite là idempotent nhưng tốn kém khi bảng lớn.}

\subsection{Lab 8.2: nạp tăng dần với high-water mark và MERGE}
Tạo bảng trạng thái \code{lake.ops.load\_state(source STRING, hwm TIMESTAMP)}. Mỗi lần chạy làm đúng thứ tự sau:
\begin{enumerate}
  \item Đọc \code{hwm} của \code{src.orders}; chốt cận trên \code{hi} bằng thời điểm bắt đầu run.
  \item Đọc từ PostgreSQL các dòng có \code{order\_purchase\_timestamp > hwm} và \code{<= hi}.
  \item Khử trùng theo \code{order\_id}, rồi \code{MERGE} vào \code{lake.staging.orders} như dưới.
  \item Chỉ sau khi \code{MERGE} commit xong mới ghi \code{hwm} mới vào \code{lake.ops.load\_state}.
\end{enumerate}
\begin{codeblock}
MERGE INTO lake.staging.orders t
USING batch_orders s
ON t.order_id = s.order_id
WHEN MATCHED AND s.updated_at > t.updated_at THEN UPDATE SET *
WHEN NOT MATCHED THEN INSERT *
\end{codeblock}
Chuyển một tháng từ \code{src.orders\_future} sang \code{src.orders}, chạy pipeline; lặp lại ba tháng. Sau đó chạy lại cùng một lần chạy hai lần, và giả lập lỗi bằng cách dừng job giữa bước 3 và bước 4.
\hoanthanh{sau mọi lần chạy lại và lần lỗi giả lập, số dòng và tổng \code{order\_id} khác nhau của \code{lake.staging.orders} khớp với \code{src.orders}; bạn giải thích được vì sao việc \code{MERGE} và ghi \code{hwm} nằm ở hai commit khác nhau vẫn an toàn.}

\subsection{Lab 8.3: dữ liệu tới muộn và dòng bị sửa}
\begin{enumerate}
  \item Chèn vào \code{src.orders} một đơn có \code{order\_purchase\_timestamp} cũ hơn \code{hwm} hiện tại (giả lập dòng tới muộn). Chạy pipeline và xác nhận nó bị bỏ sót.
  \item Đổi \code{order\_status} của 50 đơn cũ từ \code{shipped} sang \code{delivered} và đặt \code{updated\_at = now()}. Chạy pipeline và xác nhận cursor theo thời điểm mua không thấy thay đổi.
  \item Sửa pipeline: đổi cursor sang \code{updated\_at} và thêm cửa sổ nhìn lại 3 ngày (đọc từ \code{hwm - 3 days}). Chạy lại và kiểm cả hai trường hợp đã được bắt.
  \item Xoá cứng một đơn ở nguồn; ghi lại rằng không cursor nào bắt được việc này, và nêu cách bạn sẽ xử lý (đối soát toàn bộ định kỳ, xoá mềm ở nguồn, hoặc CDC).
\end{enumerate}
\hoanthanh{bạn có bảng ba cột: thay đổi ở nguồn, cursor theo \code{order\_purchase\_timestamp} có bắt được không, cursor theo \code{updated\_at} cộng nhìn lại có bắt được không.}

\subsection{Lab 8.4 (tuỳ chọn): CDC với Debezium}
\begin{enumerate}
  \item Chạy \href{https://debezium.io/documentation/reference/stable/tutorial.html}{Debezium Tutorial} theo đúng image và phiên bản trong trang.
  \item Tạo, sửa, xoá một dòng; đọc các sự kiện và chỉ ra \code{before}, \code{after}, \code{op} và tombstone.
  \item So với Lab 8.3: liệt kê thay đổi nào CDC bắt được mà cursor không bắt được.
\end{enumerate}
\hoanthanh{bạn chỉ ra được sự kiện \code{op=d} và tombstone đi sau nó cho cùng một khoá.}

\section{Kế hoạch theo tuần}
\begin{xltabular}{\textwidth}{@{}P{1.4cm} L P{1.4cm}@{}}
\toprule
\thd{tuần} & \thd{việc} & \thd{giờ} \\
\midrule
\endhead
\bottomrule
\endfoot
1 & Airbyte Sync Modes, dlt cursor và lag, Airflow Best Practices mục Creating a task, Iceberg Spark Writes, Spark JDBC; Lab 8.1 và 8.2 & 10 \\
2 & Debezium Architecture và PostgreSQL connector, Kafka Connect, Iceberg Evolution, Airflow Backfill; Lab 8.3; thử thách kỹ năng & 10 \\
thêm & Lab 8.4 với Debezium; khoá Kafka Connect trên Confluent Developer & 5 \\
\end{xltabular}

\begin{thuthach}
Xây pipeline tăng dần cho \code{order\_items} và \code{order\_payments} của Olist vào tầng staging. Hai bảng này không có cột thời gian riêng; bạn tự chọn cursor và khoá (gợi ý: khoá tổng hợp), tự quyết cách xử lý dữ liệu tới muộn và cách backfill lại một tháng bất kỳ. Pipeline phải chạy lại được nhiều lần mà không trùng dòng.
\hoanthanh{chạy toàn bộ pipeline ba lần liên tiếp và một lần backfill tháng 03/2018 cho cùng số dòng và cùng tổng \code{payment\_value} như khi tính trực tiếp trên CSV.}
\end{thuthach}

\section{Câu hỏi tự kiểm tra}
\begin{enumerate}
  \item Bạn ghi dữ liệu vào đích rồi job chết trước khi lưu cursor. Lần sau chuyện gì xảy ra, và vì sao chấp nhận được?
  \item Nếu đảo thứ tự, lưu cursor trước rồi mới ghi dữ liệu, rủi ro là gì?
  \item Kể ba loại thay đổi ở nguồn mà cursor trên \code{updated\_at} bỏ sót hoặc dễ bỏ sót.
  \item \code{INSERT INTO staging SELECT * FROM raw\_new} có idempotent không? Viết lại cho idempotent.
  \item Khi nào chọn \code{MERGE}, khi nào chọn overwrite partition động?
  \item Cửa sổ nhìn lại giải quyết vấn đề gì, và cái giá là gì?
  \item High-water mark của batch khác watermark của streaming thế nào?
  \item CDC dựa trên log cho bạn gì mà \code{SELECT ... WHERE updated\_at > ?} không cho?
  \item Vì sao \code{REPLICA IDENTITY} ảnh hưởng tới trường \code{before} trong sự kiện Debezium?
  \item Nguồn thêm một cột và đổi tên một cột khác. Tầng raw, staging và marts nên phản ứng thế nào?
\end{enumerate}

\section{Điểm dễ nhầm}
\begin{xltabular}{\textwidth}{@{}P{4.2cm} L P{3.4cm}@{}}
\toprule
\thd{điểm dễ nhầm} & \thd{thực tế} & \thd{docs chính thức} \\
\midrule
\endhead
\bottomrule
\endfoot
Nhiều người nghĩ idempotent nghĩa là nguồn không có bản trùng & Idempotent nghĩa là chạy lại cho cùng kết quả; nguồn vẫn có thể gửi lại dữ liệu, đích phải chịu được & \href{https://airflow.apache.org/docs/apache-airflow/stable/best-practices.html}{Airflow Best Practices} \\
Nạp tăng dần theo \code{updated\_at} bắt mọi thay đổi & Dòng bị xoá cứng không bao giờ xuất hiện trong truy vấn; CDC dựa trên log thấy được sự kiện xoá & \href{https://debezium.io/documentation/reference/stable/connectors/postgresql.html}{Debezium PostgreSQL} \\
Retry làm pipeline an toàn & Retry chỉ an toàn khi task idempotent; nếu không, mỗi lần retry có thể thêm bản trùng & \href{https://airflow.apache.org/docs/apache-airflow/stable/best-practices.html}{Airflow Best Practices} \\
\code{MERGE INTO} tự khử trùng dữ liệu nguồn & Nhiều dòng nguồn khớp một dòng đích thì câu lệnh báo lỗi; phải khử trùng nguồn trước & \href{https://iceberg.apache.org/docs/latest/spark-writes/}{Iceberg Spark Writes} \\
Khử trùng ở biên cursor luôn bật & Trong dlt, khử trùng ở biên tắt khi dùng \code{lag} hoặc khi đặt \code{primary\_key} rỗng & \href{https://dlthub.com/docs/general-usage/incremental/cursor}{dlt cursor} \\
Backfill chỉ là chạy lại DAG & Backfill chỉ đúng khi mỗi lần chạy ghi đè đúng partition của nó và không phụ thuộc thời điểm hiện tại & \href{https://airflow.apache.org/docs/apache-airflow/stable/best-practices.html}{Airflow Best Practices} \\
Debezium tự giữ dữ liệu cũ của dòng bị sửa & Trường \code{before} đầy đủ chỉ khi bảng đặt \code{REPLICA IDENTITY FULL} & \href{https://debezium.io/documentation/reference/stable/connectors/postgresql.html}{Debezium PostgreSQL} \\
\end{xltabular}
